\section{总结与延伸}

\subsection{核心要点回顾}

\begin{enumerate}
    \item \textbf{中间步骤}是提升 LLM 推理能力的核心机制，无论采用 training、fine-tuning 还是 prompting。
    \item \textbf{Self-Consistency}通过多次采样和多数投票，从概率角度显著提升推理准确率。
    \item \textbf{推理策略}从 few-shot CoT 到 zero-shot CoT，再到 analogical reasoning 和无提示解码，触发方式日趋灵活。
    \item \textbf{局限性}——上下文干扰、自我纠错失败、前提顺序敏感——指明了未来改进方向。
\end{enumerate}

\subsection{对 LLM Agent 的启示}

\begin{itemize}
    \item \textbf{工具调用与外部反馈}：Agent 可以通过调用代码执行器、搜索引擎等工具提供外部 Oracle，弥补 LLM 无法自我纠错的短板。
    \item \textbf{检索增强需谨慎}：RAG 引入的文档可能包含无关信息，反而干扰推理；需要精准的检索和过滤。
    \item \textbf{信息组织}：Agent 工作流中信息的呈现顺序应与推理逻辑对齐。
    \item \textbf{Scaling 推理}：Self-Consistency 和 Analogical Reasoning 可以作为 test-time compute scaling 的手段。
\end{itemize}

\subsection{拓展阅读}

\begin{itemize}
    \item Wei et al., ``Chain-of-Thought Prompting Elicits Reasoning in Large Language Models,'' NeurIPS 2022.
    \item Wang et al., ``Self-Consistency Improves Chain of Thought Reasoning in Language Models,'' ICLR 2023.
    \item Zhou et al., ``Least-to-Most Prompting Enables Complex Reasoning in Large Language Models,'' ICLR 2023.
    \item Kojima et al., ``Large Language Models are Zero-Shot Reasoners,'' NeurIPS 2022.
    \item Yasunaga et al., ``Large Language Models as Analogical Reasoners,'' ICLR 2024.
    \item Huang et al., ``Large Language Models Cannot Self-Correct Reasoning Yet,'' ICLR 2024.
    \item P\'{o}lya, ``How to Solve It,'' Princeton University Press, 1945.
\end{itemize}

%% --- Fin del contenido --- %%

\end{document}
